% ============================================================================
% APPENDIX IX: METHOD-RESEARCH - Complementarity Research Agenda
% ============================================================================
% Category: METHOD (Research Methodology)
% Version: 1.0
% Status: EXPLORATORY
% Dependencies: VII (GTC), VIII (METHOD-COMP), B (Complementarity)
% ============================================================================

\chapter{Complementarity Research Agenda: Open Questions and Future Directions}
\label{app:method-research}

% ----------------------------------------------------------------------------
% HEADER BLOCK
% ----------------------------------------------------------------------------
\begin{tcolorbox}[colback=purple!5!white,colframe=purple!75!black,title=\textbf{Appendix IX: METHOD-RESEARCH}]
\textbf{Full Title:} Complementarity Research Agenda: Open Questions and Future Directions

\textbf{Category:} METHOD (Research Methodology)

\textbf{Purpose:} Identifies, formalizes, and proposes research strategies for the major open questions in complementarity theory that extend beyond current EBF foundations.

\textbf{Status:} EXPLORATORY - This appendix maps uncharted territory.

\textbf{Relation to Existing Work:}
\begin{itemize}[noitemsep]
    \item Appendix VII (GTC): Established theory $\to$ This appendix: Extensions
    \item Appendix VIII (METHOD-COMP): Known methods $\to$ This appendix: New methods needed
\end{itemize}
\end{tcolorbox}

% ----------------------------------------------------------------------------
% CROSS-REFERENCE MAP
% ----------------------------------------------------------------------------
\begin{tcolorbox}[colback=gray!5!white,colframe=gray!75!black,title=\textbf{Cross-Reference Map}]
\small
\textbf{Builds Upon:}
\begin{itemize}[noitemsep]
    \item Appendix VII (GTC): General Theory of Complementarity
    \item Appendix VIII (METHOD-COMP): Measurement Methodology
    \item Appendix B: Complementarity Fundamentals
\end{itemize}

\textbf{Connects To:}
\begin{itemize}[noitemsep]
    \item Appendix V (MEP): Portfolio optimization extensions
    \item Appendix VI (MEQ): Equilibrium dynamics
    \item Appendix AN: LLM Monte Carlo (potential tool)
\end{itemize}

\textbf{May Generate:}
\begin{itemize}[noitemsep]
    \item Future FORMAL appendices on specific solved problems
    \item New PREDICT appendices with testable hypotheses
    \item New METHOD appendices for novel estimation techniques
\end{itemize}
\end{tcolorbox}

% ============================================================================
% INTRODUCTION
% ============================================================================
\section{Introduction: The Frontier of Complementarity Research}
\label{sec:research-intro}

\begin{tcolorbox}[colback=blue!5!white,colframe=blue!75!black,title=\textbf{Fundamental Motivation}]
The General Theory of Complementarity (GTC, Appendix VII) and the methodological toolkit (METHOD-COMP, Appendix VIII) establish a solid foundation. Yet significant questions remain \textbf{unanswered}---questions that, if solved, would substantially advance behavioral economics, intervention science, and decision theory.

This appendix serves as a \textbf{research roadmap}: formalizing what we don't know, proposing how we might find out, and identifying the theoretical and empirical tools needed.
\end{tcolorbox}

\subsection{The Six Open Questions}

We identify six major research frontiers:

\begin{table}[h]
\centering
\caption{The Six Open Questions in Complementarity Research}
\label{tab:six-questions}
\renewcommand{\arraystretch}{1.4}
\begin{tabular}{@{}clccc@{}}
\toprule
\textbf{RQ} & \textbf{Question} & \textbf{Theoretical} & \textbf{Empirical} & \textbf{Difficulty} \\
\midrule
RQ1 & Higher-Order Complementarity ($\gamma_{ijk}$) & Open & Open & $\star\star\star$ \\
RQ2 & Dynamic Complementarity ($\gamma(t)$) & Partial & Open & $\star\star$ \\
RQ3 & Endogenous Complementarity Creation & Open & Open & $\star\star\star$ \\
RQ4 & Complementarity under Uncertainty & Partial & Open & $\star\star$ \\
RQ5 & Neural Foundations of $\gamma$ & Open & Nascent & $\star\star\star$ \\
RQ6 & Causal Identification without Experiments & Active & Partial & $\star\star$ \\
\bottomrule
\end{tabular}
\end{table}

% ============================================================================
% RQ1: HIGHER-ORDER COMPLEMENTARITY
% ============================================================================
\section{RQ1: Higher-Order Complementarity}
\label{sec:rq1-higher-order}

\begin{tcolorbox}[colback=red!5!white,colframe=red!75!black,title=\textbf{Research Question 1}]
\textbf{When do THREE or more interventions exhibit synergies beyond their pairwise complementarities?}

Current theory focuses on $\gamma_{ij}$ (pairwise). But with $n$ interventions, there are potentially:
\begin{itemize}[noitemsep]
    \item $\binom{n}{2}$ pairwise terms
    \item $\binom{n}{3}$ three-way terms
    \item $\binom{n}{k}$ $k$-way terms
\end{itemize}

\textbf{Status:} Theoretically undefined, empirically untested.
\end{tcolorbox}

\subsection{Formal Definition}

\begin{definition}[Higher-Order Complementarity]
For a utility function $U: \mathbb{R}^n \to \mathbb{R}$, the \textbf{$k$-th order complementarity} among variables $\{x_{i_1}, \ldots, x_{i_k}\}$ is:
\begin{equation}
\gamma_{i_1 \ldots i_k} \equiv \frac{\partial^k U}{\partial x_{i_1} \cdots \partial x_{i_k}}
\label{eq:higher-order-gamma}
\end{equation}
\end{definition}

For $k = 3$, the \textbf{third-order complementarity} is:
\begin{equation}
\gamma_{ijk} = \frac{\partial^3 U}{\partial x_i \partial x_j \partial x_k}
\end{equation}

\subsection{Interpretation of $\gamma_{ijk}$}

\begin{tcolorbox}[colback=yellow!5!white,colframe=yellow!75!black,title=\textbf{Intuition}]
$\gamma_{ijk} > 0$ means: \textit{The complementarity between $i$ and $j$ INCREASES when $k$ is also present.}

Example: Reminders ($i$) and Incentives ($j$) might have $\gamma_{ij} = 0.1$. But when Social Proof ($k$) is added, the reminder-incentive synergy becomes stronger: $\gamma_{ij|k} = 0.2$.

This is captured by: $\gamma_{ijk} = \frac{\partial \gamma_{ij}}{\partial x_k} > 0$
\end{tcolorbox}

\subsection{The Combinatorial Challenge}

\begin{theorem}[Combinatorial Explosion]
With $n$ interventions, the number of complementarity parameters grows as:
\begin{equation}
\text{Total parameters} = \sum_{k=2}^{n} \binom{n}{k} = 2^n - n - 1
\end{equation}
\end{theorem}

\begin{example}
For $n = 10$ interventions:
\begin{itemize}[noitemsep]
    \item Pairwise terms: $\binom{10}{2} = 45$
    \item Three-way terms: $\binom{10}{3} = 120$
    \item Four-way terms: $\binom{10}{4} = 210$
    \item Total: $2^{10} - 11 = 1013$ parameters
\end{itemize}
\end{example}

\subsection{Open Problems for RQ1}

\begin{enumerate}
    \item \textbf{Sparsity Assumption:} Under what conditions can we assume most $\gamma_{ijk} \approx 0$?

    \item \textbf{Hierarchical Structure:} Does $|\gamma_{ijk}| \leq \min(|\gamma_{ij}|, |\gamma_{jk}|, |\gamma_{ik}|)$?

    \item \textbf{Sign Consistency:} If $\gamma_{ij}, \gamma_{jk}, \gamma_{ik} > 0$, must $\gamma_{ijk} > 0$?

    \item \textbf{Experimental Design:} What is the minimum sample size for a $2^k$ factorial with $k > 3$?

    \item \textbf{Regularization:} Can LASSO/Ridge methods identify significant higher-order terms?
\end{enumerate}

\subsection{Proposed Research Strategy}

\begin{tcolorbox}[colback=green!5!white,colframe=green!75!black,title=\textbf{Research Protocol for RQ1}]
\textbf{Phase 1: Theoretical Bounds}
\begin{itemize}[noitemsep]
    \item Derive conditions under which $\gamma_{ijk}$ can be bounded by pairwise terms
    \item Establish sign propagation theorems for higher-order terms
\end{itemize}

\textbf{Phase 2: Simulation Study}
\begin{itemize}[noitemsep]
    \item Generate synthetic data with known $\gamma_{ijk}$ structure
    \item Test detection power of $2^k$ factorial designs
    \item Evaluate regularization methods (LASSO, elastic net)
\end{itemize}

\textbf{Phase 3: Empirical Pilot}
\begin{itemize}[noitemsep]
    \item Design $2^3$ factorial experiment with 3 interventions
    \item Estimate $\gamma_{123}$ directly from data
    \item Compare to product of pairwise estimates
\end{itemize}
\end{tcolorbox}

% ============================================================================
% RQ2: DYNAMIC COMPLEMENTARITY
% ============================================================================
\section{RQ2: Dynamic Complementarity}
\label{sec:rq2-dynamic}

\begin{tcolorbox}[colback=red!5!white,colframe=red!75!black,title=\textbf{Research Question 2}]
\textbf{How does complementarity evolve over time? Does $\gamma(t)$ strengthen, weaken, or oscillate?}

Current theory treats $\gamma$ as static. But behavioral interventions unfold over time, and synergies may:
\begin{itemize}[noitemsep]
    \item \textbf{Build up} through repeated exposure (habit formation)
    \item \textbf{Decay} through habituation (complementarity fatigue)
    \item \textbf{Shift} as contexts change
\end{itemize}

\textbf{Status:} Partial theoretical framework, limited empirical evidence.
\end{tcolorbox}

\subsection{Formal Framework}

\begin{definition}[Time-Varying Complementarity]
Let $\gamma_{ij}(t)$ denote the complementarity between interventions $i$ and $j$ at time $t$. The \textbf{complementarity dynamics} are governed by:
\begin{equation}
\frac{d\gamma_{ij}}{dt} = f(\gamma_{ij}, x_i, x_j, \Psi, t)
\label{eq:gamma-dynamics}
\end{equation}
\end{definition}

\subsection{Candidate Dynamic Models}

\textbf{Model 2.1: Exponential Decay}
\begin{equation}
\gamma_{ij}(t) = \gamma_{ij}^{(0)} \cdot e^{-\lambda t}
\end{equation}
Interpretation: Complementarity fades with time constant $\lambda^{-1}$.

\textbf{Model 2.2: Learning/Habituation}
\begin{equation}
\frac{d\gamma_{ij}}{dt} = \alpha \cdot (x_i \cdot x_j) - \beta \cdot \gamma_{ij}
\end{equation}
Interpretation: Joint exposure builds complementarity ($\alpha$), but there's natural decay ($\beta$).

\textbf{Model 2.3: Threshold Activation}
\begin{equation}
\gamma_{ij}(t) = \gamma_{ij}^{(max)} \cdot \mathbf{1}\left[ \int_0^t (x_i \cdot x_j) \, ds > \tau \right]
\end{equation}
Interpretation: Complementarity "activates" after cumulative joint exposure exceeds threshold $\tau$.

\textbf{Model 2.4: Context-Dependent Dynamics}
\begin{equation}
\frac{d\gamma_{ij}}{dt} = \mu(\Psi(t)) \cdot [\gamma_{ij}^*(\Psi) - \gamma_{ij}(t)]
\end{equation}
Interpretation: $\gamma$ adjusts toward context-appropriate value $\gamma^*(\Psi)$ with speed $\mu(\Psi)$.

\subsection{Open Problems for RQ2}

\begin{enumerate}
    \item \textbf{Functional Form:} Which dynamic model best describes real complementarity evolution?

    \item \textbf{Time Scales:} What is the characteristic time for $\gamma$ changes? (Days? Weeks? Months?)

    \item \textbf{Asymmetry:} Does $\gamma$ build up faster than it decays, or vice versa?

    \item \textbf{Memory Effects:} Does past $\gamma$ history affect current dynamics? (Hysteresis)

    \item \textbf{Intervention Timing:} How does sequencing affect $\gamma(t)$ trajectory?
\end{enumerate}

\subsection{Proposed Research Strategy}

\begin{tcolorbox}[colback=green!5!white,colframe=green!75!black,title=\textbf{Research Protocol for RQ2}]
\textbf{Phase 1: Panel Data Analysis}
\begin{itemize}[noitemsep]
    \item Use existing longitudinal RCT data with multiple treatment arms
    \item Estimate $\gamma(t)$ at multiple time points post-intervention
    \item Test for decay, growth, or stability patterns
\end{itemize}

\textbf{Phase 2: Model Comparison}
\begin{itemize}[noitemsep]
    \item Fit Models 2.1--2.4 to panel data
    \item Compare via AIC/BIC and out-of-sample prediction
    \item Identify dominant dynamic pattern
\end{itemize}

\textbf{Phase 3: Experimental Manipulation}
\begin{itemize}[noitemsep]
    \item Design experiment varying intervention timing
    \item Measure $\gamma$ at baseline, 1 week, 1 month, 3 months
    \item Estimate time constants for complementarity dynamics
\end{itemize}
\end{tcolorbox}

% ============================================================================
% RQ3: ENDOGENOUS COMPLEMENTARITY
% ============================================================================
\section{RQ3: Endogenous Complementarity Creation}
\label{sec:rq3-endogenous}

\begin{tcolorbox}[colback=red!5!white,colframe=red!75!black,title=\textbf{Research Question 3}]
\textbf{Can interventions CREATE complementarity where none existed before?}

Current theory assumes $\gamma$ is a parameter to be estimated. But what if:
\begin{itemize}[noitemsep]
    \item Certain interventions can \textit{induce} new synergies?
    \item Training or framing can make independent actions become complements?
    \item "Complementarity engineering" is possible?
\end{itemize}

\textbf{Status:} Conceptually novel, no formal theory, no empirical tests.
\end{tcolorbox}

\subsection{Formal Framework}

\begin{definition}[Endogenous Complementarity]
Complementarity is \textbf{endogenous} if there exists an intervention $k$ such that:
\begin{equation}
\gamma_{ij}(x_k = 1) \neq \gamma_{ij}(x_k = 0)
\label{eq:endogenous-gamma}
\end{equation}
That is, the presence of $k$ changes the complementarity between $i$ and $j$.
\end{definition}

\begin{definition}[Complementarity Creation]
Intervention $k$ \textbf{creates complementarity} between $i$ and $j$ if:
\begin{equation}
\gamma_{ij}(x_k = 0) = 0 \quad \text{and} \quad \gamma_{ij}(x_k = 1) > 0
\end{equation}
\end{definition}

\begin{definition}[Complementarity Destruction]
Intervention $k$ \textbf{destroys complementarity} if:
\begin{equation}
\gamma_{ij}(x_k = 0) > 0 \quad \text{and} \quad \gamma_{ij}(x_k = 1) = 0
\end{equation}
\end{definition}

\subsection{Mechanisms of Complementarity Creation}

\begin{table}[h]
\centering
\caption{Hypothesized Mechanisms of Endogenous Complementarity}
\label{tab:endo-mechanisms}
\renewcommand{\arraystretch}{1.3}
\begin{tabular}{@{}lp{8cm}@{}}
\toprule
\textbf{Mechanism} & \textbf{Description} \\
\midrule
\textbf{Framing} & Presenting $i$ and $j$ as "part of a package" creates perceived complementarity \\
\textbf{Training} & Teaching about synergies makes people perceive/seek complementarities \\
\textbf{Platform} & A third intervention provides infrastructure for $i$-$j$ synergy \\
\textbf{Identity} & Activating an identity that values both $i$ and $j$ creates complementarity \\
\textbf{Habit Bundling} & Pairing $i$ and $j$ repeatedly conditions complementarity \\
\bottomrule
\end{tabular}
\end{table}

\subsection{The Complementarity Engineering Problem}

\begin{problem}[Optimal Complementarity Design]
Given a target behavior $Y$ and available interventions $\{I_1, \ldots, I_n\}$:
\begin{equation}
\max_{\gamma^*, k^*} \quad E[Y | \gamma^*, k^*] - C(k^*)
\end{equation}
subject to: $\gamma^*$ is achievable via complementarity-creating intervention $k^*$.
\end{problem}

\subsection{Open Problems for RQ3}

\begin{enumerate}
    \item \textbf{Existence:} Under what conditions can complementarity be created?

    \item \textbf{Persistence:} Once created, how long does endogenous $\gamma$ last?

    \item \textbf{Bounds:} Is there a maximum $\gamma$ that can be engineered?

    \item \textbf{Reversibility:} Can destroyed complementarity be restored?

    \item \textbf{Ethics:} What are the normative implications of "complementarity engineering"?
\end{enumerate}

\subsection{Proposed Research Strategy}

\begin{tcolorbox}[colback=green!5!white,colframe=green!75!black,title=\textbf{Research Protocol for RQ3}]
\textbf{Phase 1: Existence Proof}
\begin{itemize}[noitemsep]
    \item Design experiment with baseline $\gamma_{ij} \approx 0$
    \item Introduce framing intervention that presents $i$ and $j$ as package
    \item Test if post-framing $\gamma_{ij} > 0$
\end{itemize}

\textbf{Phase 2: Mechanism Identification}
\begin{itemize}[noitemsep]
    \item Test each mechanism (framing, training, platform, identity, bundling) separately
    \item Identify which mechanisms successfully create $\gamma$
    \item Estimate effect sizes for each mechanism
\end{itemize}

\textbf{Phase 3: Persistence Analysis}
\begin{itemize}[noitemsep]
    \item Track created $\gamma$ over time (weeks, months)
    \item Test if created complementarity persists without reinforcement
    \item Identify factors that sustain or erode created $\gamma$
\end{itemize}
\end{tcolorbox}

% ============================================================================
% RQ4: COMPLEMENTARITY UNDER UNCERTAINTY
% ============================================================================
\section{RQ4: Complementarity under Uncertainty}
\label{sec:rq4-uncertainty}

\begin{tcolorbox}[colback=red!5!white,colframe=red!75!black,title=\textbf{Research Question 4}]
\textbf{How does complementarity behave under risk and ambiguity?}

Real decisions involve uncertainty. Questions:
\begin{itemize}[noitemsep]
    \item Do synergies hold when outcomes are probabilistic?
    \item Is $\gamma_{\text{risky}} = \gamma_{\text{certain}}$?
    \item How does ambiguity aversion affect perceived complementarity?
\end{itemize}

\textbf{Status:} Partial theoretical framework (expected utility), limited extension to complementarity.
\end{tcolorbox}

\subsection{Formal Framework}

\begin{definition}[Complementarity under Risk]
For a lottery $L$ over outcomes with utility function $U$, the \textbf{expected complementarity} is:
\begin{equation}
\gamma_{ij}^{EU} = E\left[ \frac{\partial^2 U}{\partial x_i \partial x_j} \right] = \int \gamma_{ij}(\omega) \, dP(\omega)
\label{eq:expected-gamma}
\end{equation}
\end{definition}

\begin{theorem}[Jensen's Inequality for Complementarity]
If $U$ is concave in both $x_i$ and $x_j$, then:
\begin{equation}
\gamma_{ij}^{EU} \leq \gamma_{ij}(E[\omega])
\end{equation}
\textit{Uncertainty reduces effective complementarity.}
\end{theorem}

\subsection{Ambiguity and Complementarity}

Under Knightian uncertainty (unknown probabilities):

\begin{conjecture}[Ambiguity-Complementarity Hypothesis]
Ambiguity aversion reduces perceived complementarity:
\begin{equation}
\gamma_{ij}^{\text{ambiguity}} < \gamma_{ij}^{\text{risk}} < \gamma_{ij}^{\text{certainty}}
\end{equation}
\textit{Rationale:} Ambiguity-averse agents discount uncertain synergies more than uncertain individual effects.
\end{conjecture}

\subsection{Open Problems for RQ4}

\begin{enumerate}
    \item \textbf{Magnitude:} How much does uncertainty discount complementarity?

    \item \textbf{Asymmetry:} Is positive $\gamma$ discounted more than negative $\gamma$ under uncertainty?

    \item \textbf{Source Dependence:} Does the source of uncertainty (intervention vs. outcome) matter?

    \item \textbf{Learning:} Does uncertainty about $\gamma$ decrease with experience?

    \item \textbf{Optimal Design:} How should interventions be designed when $\gamma$ is uncertain?
\end{enumerate}

% ============================================================================
% RQ5: NEURAL FOUNDATIONS
% ============================================================================
\section{RQ5: Neural Foundations of Complementarity}
\label{sec:rq5-neural}

\begin{tcolorbox}[colback=red!5!white,colframe=red!75!black,title=\textbf{Research Question 5}]
\textbf{What is the neural basis of complementarity perception and processing?}

Questions:
\begin{itemize}[noitemsep]
    \item Is there a "synergy detection" mechanism in the brain?
    \item Do additive vs. synergistic thinkers show different neural signatures?
    \item Can neural evidence validate the $\gamma$ construct?
\end{itemize}

\textbf{Status:} No direct research on complementarity; related work on integration/bundling exists.
\end{tcolorbox}

\subsection{Hypothesized Neural Mechanisms}

\begin{table}[h]
\centering
\caption{Candidate Neural Substrates for Complementarity Processing}
\label{tab:neural-substrates}
\renewcommand{\arraystretch}{1.3}
\begin{tabular}{@{}llp{6cm}@{}}
\toprule
\textbf{Region} & \textbf{Function} & \textbf{Relation to $\gamma$} \\
\midrule
vmPFC & Value integration & May compute integrated utility with $\gamma$ \\
dlPFC & Executive control & May modulate attention to synergies \\
Insula & Risk/uncertainty & May process uncertainty about $\gamma$ \\
ACC & Conflict monitoring & May detect tension between complements/substitutes \\
Striatum & Reward prediction & May encode expected synergy value \\
\bottomrule
\end{tabular}
\end{table}

\subsection{Proposed Neuroscience Protocol}

\begin{tcolorbox}[colback=green!5!white,colframe=green!75!black,title=\textbf{fMRI Experiment Design}]
\textbf{Task:} Choice between bundles with varying complementarity

\textbf{Conditions:}
\begin{enumerate}[noitemsep]
    \item Complements ($\gamma > 0$): Items with synergy (e.g., bread + butter)
    \item Substitutes ($\gamma < 0$): Items with trade-off (e.g., Coke + Pepsi)
    \item Independent ($\gamma = 0$): Unrelated items (e.g., bread + stapler)
\end{enumerate}

\textbf{Analysis:}
\begin{itemize}[noitemsep]
    \item Parametric modulation by $\gamma$ value
    \item Contrast: Complements $>$ Independent
    \item Individual differences: Correlate neural $\gamma$-sensitivity with behavioral $\gamma$-weights
\end{itemize}
\end{tcolorbox}

\subsection{Open Problems for RQ5}

\begin{enumerate}
    \item \textbf{Localization:} Which brain regions encode $\gamma$?

    \item \textbf{Computation:} Is $\gamma$ computed explicitly or emerges from value integration?

    \item \textbf{Individual Differences:} Do some people have stronger neural $\gamma$-sensitivity?

    \item \textbf{Malleability:} Can neural $\gamma$-processing be trained or enhanced?

    \item \textbf{Disorders:} Are there conditions characterized by impaired complementarity processing?
\end{enumerate}

% ============================================================================
% RQ6: CAUSAL IDENTIFICATION
% ============================================================================
\section{RQ6: Causal Identification without Experiments}
\label{sec:rq6-identification}

\begin{tcolorbox}[colback=red!5!white,colframe=red!75!black,title=\textbf{Research Question 6}]
\textbf{When can complementarity be causally identified from observational data?}

Experiments are the gold standard for $\gamma$ estimation, but:
\begin{itemize}[noitemsep]
    \item Many intervention combinations cannot be randomized
    \item Some $\gamma$ values are only observable in natural settings
    \item We need methods for observational complementarity inference
\end{itemize}

\textbf{Status:} Active research area in econometrics; specific application to $\gamma$ limited.
\end{tcolorbox}

\subsection{The Identification Problem}

\begin{definition}[Confounded Complementarity]
Observed complementarity $\hat{\gamma}_{ij}$ may be confounded if:
\begin{equation}
\hat{\gamma}_{ij} = \gamma_{ij}^{\text{true}} + \text{Cov}(x_i, x_j | U) + \text{Selection Bias}
\end{equation}
where selection into joint treatment ($x_i = x_j = 1$) is correlated with unobserved factors.
\end{definition}

\subsection{Identification Strategies}

\textbf{Strategy 6.1: Instrumental Variables for Interactions}

\begin{theorem}[IV Identification of $\gamma$]
If instruments $Z_i, Z_j$ satisfy:
\begin{enumerate}
    \item Relevance: $\text{Cov}(Z_i, x_i) \neq 0$, $\text{Cov}(Z_j, x_j) \neq 0$
    \item Exclusion: $\text{Cov}(Z_i, \epsilon) = 0$, $\text{Cov}(Z_j, \epsilon) = 0$
    \item Interaction validity: $Z_i \cdot Z_j$ is valid for $x_i \cdot x_j$
\end{enumerate}
Then $\gamma_{ij}$ is identified via 2SLS with interaction instrument.
\end{theorem}

\textbf{Strategy 6.2: Regression Discontinuity with Interactions}

If treatments $x_i$ and $x_j$ have discontinuities at thresholds $c_i$ and $c_j$:
\begin{equation}
\gamma_{ij}^{RD} = \lim_{r \to c_i, s \to c_j} E[Y | r^+, s^+] - E[Y | r^+, s^-] - E[Y | r^-, s^+] + E[Y | r^-, s^-]
\end{equation}

\textbf{Strategy 6.3: Partial Identification / Bounds}

When point identification fails, establish bounds:
\begin{equation}
\gamma_{ij}^{LB} \leq \gamma_{ij}^{\text{true}} \leq \gamma_{ij}^{UB}
\end{equation}

\subsection{Open Problems for RQ6}

\begin{enumerate}
    \item \textbf{Instrument Validity:} How to construct valid instruments for interaction terms?

    \item \textbf{Tight Bounds:} Under what assumptions can bounds on $\gamma$ be tightened?

    \item \textbf{Machine Learning:} Can causal forests estimate heterogeneous $\gamma(X)$?

    \item \textbf{Sensitivity Analysis:} How sensitive is $\hat{\gamma}$ to unobserved confounding?

    \item \textbf{Aggregation:} How to aggregate $\gamma$ estimates from multiple observational studies?
\end{enumerate}

% ============================================================================
% SYNTHESIS: RESEARCH ROADMAP
% ============================================================================
\section{Synthesis: A Research Roadmap}
\label{sec:roadmap}

\begin{table}[h]
\centering
\caption{Prioritized Research Agenda}
\label{tab:roadmap}
\renewcommand{\arraystretch}{1.3}
\begin{tabular}{@{}clccc@{}}
\toprule
\textbf{Priority} & \textbf{Question} & \textbf{Timeline} & \textbf{Resources} & \textbf{Impact} \\
\midrule
1 & RQ6: Causal Identification & Near-term & Low & High \\
2 & RQ2: Dynamic $\gamma(t)$ & Near-term & Medium & High \\
3 & RQ1: Higher-Order $\gamma_{ijk}$ & Medium-term & High & Medium \\
4 & RQ4: $\gamma$ under Uncertainty & Medium-term & Medium & Medium \\
5 & RQ3: Endogenous $\gamma$ Creation & Long-term & High & High \\
6 & RQ5: Neural Foundations & Long-term & Very High & Medium \\
\bottomrule
\end{tabular}
\end{table}

\subsection{Quick Wins (Near-Term)}

\begin{tcolorbox}[colback=green!5!white,colframe=green!75!black]
\textbf{Recommended Starting Points:}

\begin{enumerate}
    \item \textbf{Re-analyze existing RCTs:} Many factorial experiments exist. Re-estimate with focus on $\gamma$ dynamics.

    \item \textbf{Develop IV toolkit:} Create practical guidance for finding interaction instruments.

    \item \textbf{Pilot $2^3$ experiment:} Run small-scale test of three-way complementarity.
\end{enumerate}
\end{tcolorbox}

\subsection{Theoretical Priorities}

\begin{enumerate}
    \item \textbf{Axiomatization:} Develop axiomatic foundation for dynamic and higher-order $\gamma$.

    \item \textbf{Bounds theorems:} Prove theoretical bounds relating different $\gamma$ types.

    \item \textbf{Welfare implications:} Analyze welfare effects of complementarity engineering.
\end{enumerate}

\subsection{Methodological Development}

\begin{enumerate}
    \item \textbf{LLM-augmented estimation:} Extend AN (LLM Monte Carlo) to estimate higher-order and dynamic $\gamma$.

    \item \textbf{Causal forest extension:} Develop causal forest methods for $\gamma(X)$ heterogeneity.

    \item \textbf{Bayesian framework:} Create Bayesian approach for $\gamma$ estimation with uncertainty.
\end{enumerate}

% ============================================================================
% CONNECTIONS TO EBF
% ============================================================================
\section{Connections to EBF Framework}
\label{sec:ebf-connections}

\subsection{Integration with Existing Appendices}

\begin{table}[h]
\centering
\caption{How Research Questions Connect to EBF}
\label{tab:ebf-connections}
\renewcommand{\arraystretch}{1.3}
\begin{tabular}{@{}lp{9cm}@{}}
\toprule
\textbf{RQ} & \textbf{EBF Connection} \\
\midrule
RQ1 & Extends Appendix VII (GTC): 6 $\gamma$ types $\to$ $k$-order $\gamma$ \\
RQ2 & Connects to CORE-STAGE (AW): $\gamma(t)$ varies across BCJ phases \\
RQ3 & Relates to HHH (Toolkit): Design interventions that create $\gamma$ \\
RQ4 & Links to CORE-AWARE (AU): Uncertainty affects $\gamma$ salience \\
RQ5 & Informs CORE-HOW (B): Neural basis of complementarity \\
RQ6 & Enhances VIII (METHOD-COMP): New identification strategies \\
\bottomrule
\end{tabular}
\end{table}

\subsection{Potential New Predictions}

If research questions are resolved, new falsifiable predictions emerge:

\begin{enumerate}
    \item \textbf{PREDICT-07:} Three-way interventions outperform sum of pairwise effects by factor $1 + \gamma_{ijk}$.

    \item \textbf{PREDICT-08:} Complementarity decays with half-life $t_{1/2} = \ln(2)/\lambda$ after intervention cessation.

    \item \textbf{PREDICT-09:} Framing interventions as "packages" increases measured $\gamma$ by 20-40\%.

    \item \textbf{PREDICT-10:} Ambiguity-averse individuals show 30\% lower $\gamma$-sensitivity.

    \item \textbf{PREDICT-11:} vmPFC activation correlates with $\gamma$-weighted utility computation.
\end{enumerate}

% ============================================================================
% SUMMARY
% ============================================================================
\section{Summary}
\label{sec:summary}

\begin{tcolorbox}[colback=gray!5!white,colframe=gray!75!black,title=\textbf{Summary: The Complementarity Research Frontier}]

\textbf{The Six Open Questions:}
\begin{enumerate}[noitemsep]
    \item \textbf{Higher-Order $\gamma$:} When do 3+ interventions synergize beyond pairwise?
    \item \textbf{Dynamic $\gamma(t)$:} How does complementarity evolve over time?
    \item \textbf{Endogenous $\gamma$:} Can interventions CREATE complementarity?
    \item \textbf{$\gamma$ under Uncertainty:} How does risk/ambiguity affect synergies?
    \item \textbf{Neural $\gamma$:} What is the brain basis of complementarity?
    \item \textbf{Identification:} When can $\gamma$ be estimated without experiments?
\end{enumerate}

\textbf{Near-Term Priorities:}
\begin{itemize}[noitemsep]
    \item Develop IV methods for interaction identification
    \item Analyze $\gamma$ dynamics in existing panel data
    \item Pilot $2^3$ factorial for three-way effects
\end{itemize}

\textbf{Long-Term Vision:}
A complete theory of complementarity that includes:
\begin{itemize}[noitemsep]
    \item Arbitrary-order interactions
    \item Full temporal dynamics
    \item Endogenous creation/destruction
    \item Neural grounding
    \item Robust causal identification
\end{itemize}

\end{tcolorbox}

% ============================================================================
% GLOSSARY
% ============================================================================
\section{Glossary}
\label{sec:glossary}

\begin{description}
    \item[Higher-Order Complementarity] $\gamma_{i_1 \ldots i_k}$: Interaction among $k > 2$ variables
    \item[Dynamic Complementarity] $\gamma(t)$: Time-varying complementarity
    \item[Endogenous Complementarity] $\gamma$ that can be created or destroyed by interventions
    \item[Complementarity Engineering] Deliberate design of interventions to create synergies
    \item[Interaction Instrument] IV for the product term $x_i \cdot x_j$
    \item[Complementarity Fatigue] Decay of synergies with repeated exposure
\end{description}

% ============================================================================
% REFERENCES
% ============================================================================
\section{References}
\label{sec:references}

\nocite{bcm_master}

This appendix draws on:
\begin{itemize}[noitemsep]
    \item Appendix VII (GTC): Established complementarity theory
    \item Appendix VIII (METHOD-COMP): 128 methodological papers
    \item Emerging literature on causal inference, neural economics, dynamic models
\end{itemize}

\textbf{Key papers for future research:}
\begin{enumerate}
    \item Higher-order interactions: Stein et al. (2021) on higher-order statistics
    \item Dynamic treatment effects: Callaway \& Sant'Anna (2021)
    \item Causal forests for interactions: Athey \& Wager (2019)
    \item Neural value integration: Bartra et al. (2013)
\end{enumerate}

% ============================================================================
% OPEN ISSUES
% ============================================================================
\section{Open Issues}
\label{sec:open-issues}

\begin{enumerate}
    \item \textbf{Prioritization:} Which research question should the community tackle first?

    \item \textbf{Collaboration:} How to coordinate multi-disciplinary research (economics, psychology, neuroscience)?

    \item \textbf{Data Infrastructure:} What shared datasets would accelerate $\gamma$ research?

    \item \textbf{Replication:} How to ensure robustness of complementarity findings?

    \item \textbf{Translation:} How to communicate complementarity research to practitioners?
\end{enumerate}

\end{chapter}
